\chapter{Discussion and conclusions \label{chapter:conc}} 

The modelling paradigm of \gls{abm} has held out the promise of freedom from all earlier model-building restrictions, able to create `virtual worlds' (section \ref{sec:virtualworlds}) with objects mapping perfectly to real-world structures. If this promise were straightforwardly true, \gls{abm} would seem perfectly suited to the `big questions' of spatial economics: why does spatial economic activity happen in the way it does? How does it lead to the economic landscape seen in reality? How does it change? How do these all happen in a spatial economy made up of agents separated by time and space making individual, but ultimately massively interconnected, decisions?

This conclusion argues that any attempt to answer those questions using \gls{abm} must take the limitations of \gls{abm} more seriously. This is not a negative conclusion: it reveals a great deal of scope for further work usually hidden under layers of complexity. \gls{abm} has huge promise as a tool for digging into the most difficult spatial economic problems; understanding its limitations is a vital step in living up to that promise.

A set of specific dependencies and assumptions have been revealed through the model development process; the upcoming discussion uses these as the focus for highlighting both the limitations of \gls{abm} and new opportunities these limitations reveal. 

The first aim of the thesis was to:

\begin{itemize}
\item \textbf{Develop an agent-based spatial economic model framework in which agents with heterogenous locations can optimise welfare given changing costs}
\end{itemize}

This has the following objectives connected to it:

\begin{itemize}
	\item Develop a spatial agent model framework and explain its workings in detail in the thesis;
	\item Compare agent-based and economic approaches to model-building (with specific emphasis on \gls{ge}) and use this analysis to inform the thesis' model development.
\end{itemize}

Chapters \ref{chapter:modelfw} through \ref{chapter:connectedModels} presented the framework developed to do this and key results along the way. These described the route to producing the transmission belt model, which successfully enables agents with heterogenous locations to optimise welfare given changing costs, using `laws of motion' (section \ref{sec:ace}) plus damping, allowing Firms to target a stock level and People to maximise utility. The feedback between these two drives a transition between production equilibria where one or more efficient producers emerge.
%This thesis aimed for a `general' agent-based spatial economic model, narrowly defined as able to link supply and demand in a spatial setting where agents, making separate economic decisions, could produce a stable equilibrium. 

As well as a series of testing models, chapter \ref{chapter:models} focused on two modelling scenarios, each examining a separate issue. The `mobile People' models (section \ref{sec:results}) looked at spatial equilibrium by using a `density cost' proxy that avoided coordination issues. The `mobile Firms' models (section \ref{sec:mobFirms}) concentrated on the coordination issue, examining how decisions in a simple two-Firm scenario interact. 

\gls{ge} has been examined in detail: both the `nuts and bolts' of its modelling approach and its overall modelling philosophy (as well as that of economics more broadly). This thesis has taken seriously the arguments about spatial economics put forward by \gls{ge}. In particular, the need for a link between supply (through some theorisation of production) and demand has been at the heart of the modelling approach. 

The benefits of closely examining \gls{ge} have been two-fold. Firstly, many concrete ideas have been used, not least the \gls{ces} function throughout the modelling framework. Secondly, and just as importantly, it has revealed a set of challenging problems outside the norm of concerns for \gls{abm}. Where the default \gls{abm} question has tended to be `Hayekian' (section \ref{sec:ace}) (`how can it come about through decentralised actions?'), keeping the focus on \gls{ge} concerns about spatial supply and demand has pushed the thesis model framework to find a route that works in an \gls{abm} setting.

Doing this has revealed a clear set of dependencies and assumptions. The result is far from negative, as this final chapter hopes to show. Rather than striving for an \gls{abm} ideal that supersedes all previous modelling paradigms, turning round and looking the other way, towards the tangle of dependencies and assumptions, reveals a rich seam of possibilities and allows a more open comparison to other, equally valid, modelling approaches.

The second aim of the thesis is:

\begin{itemize}
\item \textbf{Use the development process of the model framework to answer the following: what obstacles are in the way of \gls{abm} exploiting its potential as a tool for spatial economic analysis?}
\end{itemize}

The following objectives are connected to the second aim:

\begin{itemize}
	\item Identify obstacles revealed in modelling development at the coding level and a route through them;
	\item Identify obstacles revealed by the comparison of different approaches to model-building in \gls{abm} and economics.
\end{itemize}
%Where the first aim is about finding a successful path through to a model linking spatial supply and demand

These two objectives are interlinked. Solving problems in the process of building the thesis models meant identifying and overcoming a set of obstacles. In combination with examining the theory of modelling, this process has been used to think about obstacles for agent-based spatial economics and why \gls{abm} seems to have avoided the `big questions' of spatial economics.

Obstacles discovered during the development process have ranged from highly specific coding choices through to broad theoretical ideas. An example coding choice would be working out the need for a `smooth' \gls{eos} function rather than the \gls{coremodel}'s fixed/marginal function (sections \ref{sec:prodfuncts} and \ref{sec:abmprod}). The effect of considering theory has been broader. Examining the roots of \gls{abm} and \gls{ge} has been, to paraphrase Hoover (section \ref{sec:mot}) `a quotidian practice, not abstract epistemology', directly applicable at various model turning points. The main effect has been to provide  `permission' to explore a space of modelling choices not usually considered `valid' in \gls{abm} - not least freely using ideas from neoclassical economics, given that many \gls{abm} theorists believe the field to be anachronistic.
%The nature of the model building endeavour is a `praxis'; it is difficult to predict what obstacles may arise without actually building a model. 

Through this process of developing the thesis' \gls{abm} spatial economic framework, the following two obstacles have been identified:

\begin{enumerate}
	\item A lack of research directly tackling the space of dependencies in \gls{abm} model development. Rather than studying dependencies abstractly, further work should naturally result from model-building while carefully considering the algorithmic `nuts and bolts' of the process. This is a promising route for \gls{abm} \emph{because} of its flexibility: the possibilities for trying out new code structures are huge. \gls{ge} must rely on the `engine' of the \gls{dxs}; this ties it to a specific set of dependencies that, for example, rule out considering a situation where firms make decisions over a much longer timescale than people. \gls{abm} would in principle have no problem with this. 
	\item A broader sense that \gls{abm} has broken all ties with history and can learn nothing from past modelling efforts. By exploring economic ideas, the thesis suggests that in the endeavour of modelling economic systems, pre-computational thinking is still highly relevant. As the introduction discussed, while computation \emp{has} radically altered the way modelling is done, it has not changed its fundamental nature. 
\end{enumerate}

The rest of this chapter aims to illustrate the fruitfulness of taking dependencies seriously, as well as treating \gls{abm} as fully part of a much longer history of modelling thought. This ties together the two main components of the preceding thesis: the presentation of modelling choices, ideas and theory in chapters \ref{chapter:location} and \ref{chapter:ABM} and the thesis model framework results presented in chapters \ref{chapter:models} and \ref{chapter:connectedModels}. As Hoover points out (section \ref{sec:code_vs_function}), models cannot be `cleanly idealised' if they are to be `operable or realisable in a manipulable form'. Dependencies are unavoidable. Being explicit about this helps track down what they look like and what possible new routes may be open for exploration. The rest of the conclusion looks in detail at these assumptions and dependencies, in the following sections:

\begin{itemize}
	\item A direct comparison between the \gls{coremodel} and the thesis model in order to highlight some of the key assumptions and differences;
	\item Model framework dependencies identified during the building process;
	\item The different types of agent interaction the framework has examined and the assumptions involved. These are tied back to considering obstacles facing a spatial economic \gls{abm};
	\item A more detailed examination of assumptions in both the model framework and in model development more generally;
	\item The conclusion ends with a discussion of possible further work and some final comments.
\end{itemize}

\section{Comparison with the core model \label{cmcomp}}

While the thesis model framework is not meant to be a direct comparator to the \gls{coremodel}, it is informative to compare them on a number of categories to bring out the differences in both outcome and modelling philosophy of \gls{ge} and \gls{abm}, as well as setting up discussion of model dependencies in section \ref{sec:deps}. Fowler found during the process of developing his agent model that analytic choices made to gain tractability can ``actually make the problem less tractable'' in an agent setting (\citealt{fowler_taking_2007} p.283). This has certainly been the experience throughout the process of model development; this section highlights where many of these differences arose.

Di Paolo \emph{et al.} make the point that no one methodology is inherently superior (section \ref{sec:virtualworlds}), but each does have different strengths and weaknesses. For example, the approach used by \gls{ge} versus the agent models of the thesis have quite different areas where they manage to be simple and where they fall into complication (rather than complexity). The detailed presentation of the \gls{coremodel} in its original, simplest possible incarnation, illustrates plainly that it is far from simple in its components and the web of assumptions required for it to function, even while the dynamics it reveals are simple to understand. This thesis' agent approach offers a simple way to think about agents, using quite basic functions to make decisions over distance. But the process of making them interact is not simple.

From a top-level view, the difference might be described as follows. The \gls{coremodel} creates a tension between worker's income based on a wage, which is worked out from a price index, which is built on a complex market structure argument. The final thesis model has these features: People have time to spend. They exchange it for goods over a distance. Their time goes into a range of Firms each with an \gls{eos} function. Firms adjust price to keep the market cleared. Doing this allows efficient producers to emerge through cumulative causation. The core model is built on a scaffold: simplifications and useful assumptions that allow the model to build up to its final structure, losing the complexities it used to get there. The thesis models are an attempt to build on the strengths of \gls{abm} by keeping the key interactions present.

\subsubsection{Timing structure}

The \gls{coremodel} is broken into a short and long run, both with distinct dynamics - though the long run is by far the simpler of the two. Each is mathematically consistent but the bridge between them is essentially arbitrary, as long as workers move where the real wage is highest. The short-run component of the \gls{coremodel} uses an iterative process to find its equilibrium values, but this is simply a method of finding a stable set of values for the region equations; it has no economic meaning.

The thesis models run over many time steps. This fits with Hicks' description of supply and demand functioning ``in a recursive manner'' (\citealt{hicks_market_1989} p.11). This contrast to the \gls{coremodel} shows up well in both the `mobile People' models and the transmission belt model. The `mobile People' model's use of the density cost (section \ref{chapter:models}) sees agents producing emergent morphological outcomes over model iterations with every single Person's choice having external effects. It is this recursive, external-effect interaction that produces this model's spatial equilibrium outcomes.

Similarly, the transmission belt model relies on step-by-step feedbacks to transition between production equilibria, utilising the smooth \gls{eos} function to do so. It does not produce the migration outcomes of the \gls{coremodel}; cumulative causation occurs through People deciding where to exchange their time, and this translating into production equilibria (see below for more on the production structure).

\subsubsection{Use of agents and the nature of firms and people}

In the \gls{coremodel}, the number of agents actually represented by specific mathematical objects or code are small: as section \ref{sec:fowler} says, each region has a representation of manufacturing and agricultural workers, making four `agents' in total (though these reduce to summed region income  equations). These are deemed `\gls{repag}s', standing in for all economic actors. Demand from each of these types of agent (manufacturing and agricultural) uses the Cobb Douglas' ability to hermetically seal them into separate compartments; this is very useful for the \gls{coremodel} as it allows the manufacturing side of the market to be treated separately. For an agent-based model, there may be situations where such a clean separation could be useful, but not if agents are required to have interactive demand.
%A similar contrast exists as regards the Cobb Douglas utility function, used as the top-level utility in the \gls{coremodel} for manufacturing and agricultural goods. As section \ref{sec:cobbdouglasresults} discussed, 

The number of firms in the \gls{coremodel} can be represented by a continuous variable, deduced from the models' equilibrium structure, though the number itself plays no direct role in outcomes. The number is `assumed to be large' in order for some key macroeconomic simplifications to be made, allowing vital elements of the mathematical argument to be set to constants. This is a feature of the \gls{dxs} approach the \gls{coremodel} uses to good effect, building on it to produce its whole equilibrium structure. It is presumed that when workers move to find higher utility, firms follow. The equilibrium number of firms also follows instantaneously. Firm entry and exit is able to adjust at least as fast as workers are able to migrate.
%Firms' presence is implicit in the \gls{coremodel} in that they have no direct part to play in the feedback processes involved. 

Firm and Person agents represent each side of supply and demand. While in theory, each have encapsulated decision-making processes, in practice, the structure of these decision processes is determined by the model type and highly dependent on given model choices; this point is discussed in detail in section \ref{sec:deps}. In the `mobile People' models, Firms are static sites that provide goods at fixed prices to People with fixed wages. Many People agents optimise their locations, trading off distance costs, base costs and an incurred `density cost' that allows a spatial morphology to emerge. The `mobile Firms' model uses only two Firms, competing over location on a line of People. The medium of competition (location) and the number of Firms are kept as basic as possible to illustrate a key issue of signalling in the model framework (using \gls{dmr}).

All of these, including the transmission belt model, use a relatively small number of agents (maximum $500$). Agent modelling relies on the agents having some kind of market signal to work with. This does not mean that `very large market' arguments may not be useful, but if reaction to market signals is being investigated, this becomes problematic the larger the number of agents and the noisier the environment. The relatively small number of Firms in the thesis models allows external interaction between price-setting behaviour, but is not `oligopolistic' as Firms do not directly compete.

%This is one of those points in the model where one has to ask what simplifications are justified. There are a couple of arguments to make. First, the basic one about useful but silly assumptions, and the fact that FKV recognise these and spend time addressing: ``some of its implications turn out to be sensitive to assumptions one would not want to defend.'' p.76. 
%This specific one: timing of firms etc. It is "the engine" so not up for negotiation. The ones they're discussing are peripheral to that.

\subsubsection{Production structure}

The \gls{coremodel}'s use of a fixed and marginal labour production structure is the basis for its argument about how firms enter and exit the market. The mathematical simplifications involved allow the model to remove any elasticity of substitution between goods: the production scale itself in each firm stays fixed and changes in market structure happen \emph{solely} through the implied entry and exit of firms. It is a key simplification that allows the \gls{dxs} `engine' to function.

The transmission belt model works through changing production scale. It does not allow firm entry or exit, but each Firm agent is able to use price signals to target an individually `efficient' level of stock production. At certain values for \gls{lov}, \gls{eos} and distance, the model transitions between production equilibria, where one or more Firms emerge as sites of much more efficient production.

%it presents quite specific timing and decision problems that using a `smooth' production function does not. 
Fixed and marginal requirements for production make sense in an algebraic setting, where a variable for each can play a role. In the \gls{coremodel}, the marginal requirement is used both in the market structure argument and in key normalisations allowing for much simpler regional equations. Getting this production function to work in an agent model is not necessarily an insurmountable problem, but as section \ref{sec:blindvsrational} explored, creating a functioning model with \gls{eos} as a feature presents problems of strategy and timing that must be overcome. In the thesis, this is done by assuming People leave their time with Firms between decision points, thus allowing a flow of production to take place. 

%The production function in the transmission belt model was used because it fitted well with the other model dependencies, as well as the iterative nature of this kind of modelling. 
While the smooth production function was chosen because it fitted with this `leaving time' approach to producing \gls{eos}, it also transpired it could be analysed in spreadsheet form to identify the central driving dynamic responsible for producing shifts between stable production equilibria (section \ref{sec:lov_vs_eos}; see below also).
%This process of cumulative causation can be boiled down to a simple iterative set of spreadsheet cells. In the thesis model framework, Firms and People are also able to use the same basic price-targeting mechanism to produce production equilibria within a large variety of spatial settings. 


%This is a key contrast to both production structure and market structure in the \gls{coremodel}, where cumulative causation is determined \emph{after} the short-run equilibrium is found in a quite different model `compartment'. 
Cumulative causation (or feedback) have a role in both the \gls{coremodel} and the transmission belt model. It is present in two ways in the \gls{coremodel}. Firstly, though an argument about the feedback between firms and workers. Since firms follow workers, there will always be more firms where more workers are; through \gls{lov} (played out via changes in firm number) this locks in a higher real wage (\citealt{brakman_new_2009} p.127). This creates the short-run difference in wages between regions. In the long-run, feedback of migrating workers is argued to occur. Section \ref{sec:cmdynamics} showed how this dynamic is backed up by reference to the polarity of short-run real wage differences for any given migration change. In short, migration is guaranteed to occur until a stable equilibrium is reached; this leads to either a spread of workers or the emergence of a core in one of the regions. 

Cumulative causation in the transmission belt model is due to a feedback between the difference in demand and production. A minimum of two Firms are required for this dynamic (a single Firm would get all market demand and thus its production scale is fixed). A two-Firm model can begin in symmetric equilibrium with both producing at the same scale. If demand passes a substitutability threshold (controlled by a balance between \gls{lov}, \gls{eos} and, if present, distance costs), the magnitude of the difference between demand and production at a given price can keep both Firms away from their stock target, causing a feedback that ends in a new production equilibrium with the two Firms producing at radically different scales.

The ability to develop non-symmetric demand structures is a key contrast to the \gls{coremodel}, in which goods are guaranteed to be demanded symmetrically. No changes in production structure at the firm level are possible; increasing returns occur through firm entry and exit only. The transmission belt model has the opposite dynamic: the number of Firms is fixed and production structure changes play out through production-level \gls{eos}. In section \ref{sec:twofirmsperregion}, for instance, with two Firms per region, four permutations of production structure emerge for different values of $\rho$ and distance cost.

This method, with the addition of the `avoid \gls{eos} strategy' code structure (section \ref{sec:abmprod}), allows for more heterogenous agent models with the same dynamic to reach stable equilibria, although stability is dependent on parameter settings and is affected by an interaction of \gls{lov} and the damping response time of Firms targeting a stock level.

Relating back to the `ecological criticism' discussion in section \ref{sec:ecocrit}, the ability to explicitly separate \gls{eos} from \gls{lov} potentially provides a way to examine spatial economic issues involving `Weberian' weights and density. In combination with the ability to produce non-symmetric production structures, the transmission belt model may be suitable for considering production networks; this point is returned to below in relation to dependencies, as well as in section \ref{sec:theend} discussing further work. 

\subsubsection{Market structure}

Section \ref{sec:spacesituate} outlined how market structure has determined the divergent paths of non-spatial and spatial economics, up until \gls{ge} attempted to bridge the gap. `Market structure' determines the way production and consumption connect. In the thesis model framework, this has been labelled a `transmission belt' to emphasise the explicit feedback between production output and demand. 

At root, production and consumption are conceptualised in much the same way in both the \gls{coremodel} and the thesis models, in that simple output and optimisation functions are used. The thesis has deliberately avoided any more complex structures that \gls{abm} may in theory make possible. The two approaches link production and consumption very differently, however. 

In the \gls{coremodel}, the equivalent of the `transmission belt' between production and consumption happens in a mathematical argument leading to the deduction of regional real wages. Through a series of simplifications, key equilibrium points are found. These are where the argument can state: \emph{if} supply equals demand \emph{then} the equations can be transformed thus. This happens in two places in the \gls{coremodel} leading up to the short-run equilibrium. Firstly, an optimal, demanded quantity of good (deduced through assuming some key macroeconomic quantities are constant) is plugged into firms' profit equation, which is assumed to be maximised. This is plugged into a regional price index describing a proxy for the real wage, or `the cost of a unit of utility'. Second, equations for the optimal amount of demanded good and `the optimal amount produced in equilibrium' are combined. Again, the next step is a series of re-arrangements to arrive at an equation in terms of the regional nominal wage. This is combined with the price index to produce a regional real wage.

While the presence of firms is implicit in the \gls{coremodel}, the structures it uses develop an argument for a \gls{monocomp} market structure. At equilibrium, the number of firms (and thus varieties) change; this is responsible for changes in the real wage and thus ultimately determines where workers in the model go. 

In the thesis, `partial' models examine first People's movement, then Firms' movement. Both models are driven by People's demand. There is only direct competition in these two when two Firms compete over location. The transmission belt model makes the `general' link between supply and demand by using Firms that target zero stock, where stock is produced using a smooth \gls{eos} function. The resulting equilibrium comes about through a robust link between price signals, production and demand, but introducing either a profit motive or an ability to genuinely compete has proved difficult (section \ref{sec:markup}). The \gls{coremodel} (and other static equilibrium models) avoid this problem by assuming competition has already taken place. 

A more fundamental question: what does finding this simple `transmission belt' example, using a smooth \gls{eos} function and a separate \gls{lov} function, say about increasing returns, if anything? In the simplest spreadsheet example, there is no space, two Firms and one Person. \Gls{lov} alone can determine whether one of those Firms emerges as the more-efficient producer. This happens without any of the assumptions involved in the \gls{coremodel}'s complex, murky and in places somewhat odd market structure argument. 

However, without an explicit competition dynamic, it is difficult to directly compare to a \gls{ge} approach where, in the absence of transport costs, ``free trade ensures that the prices of all varieties are equalized between countries'' (\citealt{brakman_new_2009} p.130). As section \ref{sec:markup} discussed, it may even be impossible for the transmission belt model to produce any stable competition equilibrium in its present form. The `transmission belt' does, at least, show that it is possible to produce an increasing returns dynamic with some very simple assumptions, and that both production-level \gls{eos} and \gls{lov} can interact in a feedback process to cause it. 

There is a market structure in the transmission belt model, then, but not a fully competitive one. The final section returns to this point as one of the major areas for further work.

\section{Dependencies \label{sec:deps}}

\subsection{Overview}

The concept of dependencies provides both a useful way to think about how model components connect and draws attention to a possible flaw in much of the way \gls{abm} is pursued: it is too reliant on code-mapping (section \ref{sec:mapping}) and thus to the ideal of `loosely coupled' agents interacting quite separately in their environment. This section summarises the dependencies found in the process of developing the results in chapters \ref{chapter:models} and \ref{chapter:connectedModels}. The introduction suggested that the \gls{oop} language of dependencies was useful not only for discussing agent models, but for thinking about how model assumptions in general link together or cause certain approaches to be ruled out. In this light, the original spatial market structure problems of increasing returns and distributed market activity have determined the dependencies of economic theory (section \ref {sec:spacesituate}).

Dependencies are entirely internal to a model: they create `tightly coupled' structures that dictate the form of a model's assumptions. So, for instance, pre-\gls{ge} general equilibrium models using only linear functions create a dependency requiring perfect competition. This in turn means it cannot deal with the fundamental spatial economic requirement for increasing returns (section \ref{sec:spacesituate}).

Examining dependencies also produces some very positive outcomes: it helps clearly reveal other possible modelling avenues. This is aided by the transparency and simplicity of the economics concepts used; dependencies can be hard to see under layers of model complication. Modelling avenues identified in this thesis point in a rather different direction to a lot of existing \gls{abm} work, toward `messy' coding issues where meaning and assumptions cannot be `cleanly idealised'. This is, however, still  an area where \gls{abm}'s flexibility offers a great advantage: a huge abundance of possible ways to tackle any modelling problem. The choices leading to the transmission belt model presented many branches; some of these specific dependencies are discussed below. 

%The goal for the thesis model framework was to produce a spatial supply-demand equilibrium with agents able to making independent, optimal choices from different locations. The specific goal (agents making separate spatial economic decisions based on their choice set) does not necessarily translate into a coding reality as clean as `build independent agents, allowing them to interact and observing the result.' In the process of making a model that manages to achieve this, 
\subsection{Model framework dependencies \label{sec:depsframework}}

A key determinant of the models' dependencies is the decision to use a \gls{ces} function for utility. Two parallel methods were tested (section \ref{sec:greedyvsconstrained}). First, a constrained optimisation that umbilically links good prices to People's demand. A given set of prices, through the constrained utility equations, tell People what amount of each good is optimal. For the model to link demand to supply, that demand needs to actually be met by Firms. Second, a `hill-climber' approach which converges on the mathematically perfect constrained optima the more iterations it is given, starting from a random guess at a mix of good amounts and `climbing'. 

Both of these utility methods are capable of optimising a given set of goods using a \gls{ces} function. The `hill-climber' is more flexible, but this is paid for: it is computationally expensive. It is, however, potentially able to allow People to deal with scenarios involving a lack of optimal stock amount. Zero stock is no problem for either utility approach: a good is simply not entered into the utility function. But the hill-climber is able to re-allocate any spending from goods that have partially enough stock, so it can be spent on others where more stock is available - thus deviating from the `objective' amount demanded in the constrained utility function.

Using a constrained utility function means People assessing their utility `now' and thus needing their demand met at the same point in model time. This causes a dependency in the way Firms must be structured: they cannot be given their own `instant' assessment procedure. This would be possible, for example, if each time a Firm acted, it was guaranteed to be able to clear its stock. It would then be up to People agents to adjust their price-setting over time instead. Both cannot be instantaneous. There are workarounds, but these involve `time-machine' firms, able to reset the world and re-run each action to compare different outcomes. This is possible, in terms of coding: state can be kept for the whole model and each Firm given the chance to rerun each action, with stochastics dictating the range of other Firms' responses, in order to `climb' to the best action. But this is a long way from the goal of allowing agents to make independent decisions: this gives them `spooky' knowledge of all the paths a market could take. 

If Firms were equipped to make optimal decisions `now', the hill-climber approach may also offer another route out of a key dependency that results from using the constrained utility function: the hill-climber iterations could be interpreted as model timesteps, suggesting a way for good buying to be something other than optimal `now'. Random good `slices' could continually adapt as supply changed over time without ever reaching a `perfect' optimum. The constrained utility approach could not do this.

Price competition over good cost and wages does not work without some limit on stock. If stock is potentially infinite, Firms have no restriction on the amount they can sell. If Firms attempt to optimise price as they do with location in section \ref{sec:mobFirms}, prices repeatedly bounce along next to zero or wages enter a spiral as Firms attempt to out-compete each other. Introducing a stock limit, however, introduces new complications: the model is now moving from `partial' to `general' territory. It is difficult to find any compromise between these two, as the problem of unlimited stock illustrates. For the introduction of stock to work effectively, coordination needs to take place between wages, good prices and demand. Attempting to solve this by allowing Firms a `sonar' feedback between timesteps finds itself beset by \gls{dmr} problems. Section \ref{sec:mobFirms} explored this issue by allowing two Firms to attempt optimising their location only, along a line of People with fixed, equally distanced positions. The ability of Firms to correctly respond to market signals was shown to be affected by other events within that market: it was possible for them to interpret meaning incorrectly and thus make inapproriate responses.

The exact way that costs are actually communicated is also key. In an agent economic model, what do agents pay with? Firms buying labour and selling goods require some medium of exchange. As section \ref{sec:comm_costs} discussed, many agents using some cost signalling currency may find themselves producing unintended entropy-like dynamics where the currency becomes the limiting factor and distracts from the analysis. While other cost-communication methods may be possible, the solution in this thesis has been to make People's time the currency and to subsume wage and goods into a single exchange, thus massively reducing the complexity. But that does rule out considering labour and good markets separately: another dependency.

The transmission belt model introduces its own dependencies. As section \ref{sec:abmprod} explained, `\gls{eos}' dynamics have the potential to be a thorny coordination problem of an `El Farol bar' type, caused by many agents choosing what to buy and where to work. This is solved through two assumptions. First, People are guaranteed demand is met at the price they find (as well as combining the wage/good decision into one). Second, they `leave' their time with the Firms they buy from until their next decision, when they remove it and decide anew. This allows a flow of production to take place that assumes People invest time between decision points. \Gls{eos} thus become possible and, as a result, stock targeting also becomes possible. This is a particularly tightly bound set of dependencies that the transmission belt model needs to work in a many-agent situation, though the spreadsheet example (section \ref{sec:lov_vs_eos}) illustrates how its dynamic functions without these complications.

Where the \gls{dmr} two-Firm model enabled Firms to optimise location (albeit through some noise), the price volatility inherent in the transmission belt model makes \gls{dmr} much more problematic. As section \ref{sec:partialconcs} discussed, there are workarounds that involve making sure Firms can make decisions in an environment where the signal-to-noise ratio is workable for them (which only requires them to take correct actions more often than incorrect ones on average). Given the volality of the transmission belt model's stock targeting mechanism, this would require oscillations to damp enough between actions for any action to produce a readable `sonar'. However, this would entail the entire model pausing and Firms `telling' it when to restart. While the current structure is some way from the ideal of agents making perfectly independent decisions, this approach would firmly rule that out.

This noise issue rules out even Firm location optimisation which is, as section \ref{sec:mobFirms} discussed, in principle no different from the process of searching other price spaces. Equally, as discussed in section \ref{sec:markup}, Firms' ability to optimise their profit must also rely on knowing if their own action (changing the markup target) has led to an increase in revenue.

By subsuming exchanges into one transaction, the transmission belt model builds in two other dependencies. One has been mentioned: there is no way to model a separate labour market, which may be key to allowing a stable competition dynamic, since at present People have no `sanction' against Firms skimming very nearly all production output. Relatedly, it gives the wage-labour exchange a \gls{lov} dynamic: People buy a mix of goods and in doing so also give a mix of time to Firms. Tendentiously, there are real-world parallels to be made to people working for more than one employer or in managing `portfolios' of food production sites (\citealt{rose_who_1993}), but this still nevertheless places a restriction on the model.

As mentioned above, the size of the market also matters for price signals, especially for those of the \gls{dmr} type where noise is an issue. In an agent model, agents need to be able to receive \emph{some} information to have a basis for action. Where the \gls{coremodel} \emph{relies} on an argument that the `market is large' enough to avoid any direct price interaction, in the thesis model framework, a market this large would make signals increasingly difficult for agents to use. 

The transmission belt model represents one possible path through these dependencies to a working model. Highlighting dependencies in this way suggests other modelling ideas to test, though many may involve radically different root modelling assumptions. The final section of this chapter discusses some possible further research ideas based on this.

\section{Different types of interaction \label{sec:coordoverview}}

Agent interaction is the main motivation for using \gls{abm}. For `interaction' to take place, outcomes need to not be simply ergodic: the presence of multiple agents must be a factor, not just producing stochastic outcomes where, for example, a model with one agent run many times produces the same result as many agents run once. The dependencies described in section \ref{sec:deps} underpin four distinct types of interaction in the thesis' models. This section examines how and why they differ in their source and effect, as well as discussing their relationship to assumptions.

Two interaction types avoid coordination: the density cost and the approach to \gls{eos}. Two others - Firms' \gls{dmr} location feedback and the transmission belt model's stock targeting - require some form of coordination to occur, meaning that a set of interdependent decisions determine model outcomes.
%Assumptions can be used that can circumvent coordination problems: the density cost does this. 

The \gls{eos} dynamics of the transmission belt model, as mentioned in the previous section, avoid coordination problems by allowing demand to be met while providing a method for production to change scale as the `flow' of time input changes. It is, technically, a point of interaction: the change of production scale is essential for the transmission belt mechanism and this comes about through many agents moving more demand to a subset of Firms. It is the most straightforward of the interactions discussed here. It illustrates that there are effective options for avoiding potentially insurmountable coordination problems.

The density cost provides a way to think about how individual agent actions in space can aggregate, through external effects, into spatial morphology. Disadvantages arise if the target of interest is more specific land market dynamics; the density cost has no market exchange. But this is also its advantage: People responding to it are not capable of making a `wrong' decision, as the \gls{dmr} Firms are. Whatever outcome they choose is guaranteed optimal \emph{at their point of decision}. The emergent outcome, which may be sub-optimal for some, is externalised to the landscape - it does not happen inside the agents themselves but through an accretion of location decisions and external effects. As section \ref{sec:notutilityeq} says, the result shows a path-dependent outcome: a spatial equilibrium will result, but some agents may have hold of a better location than others. On each agents' turn, however, if no better location is available, no change happens: a Nash equilibrium has been reached. 

At the other extreme, the \gls{dmr} scenario in section \ref{sec:firmsLoV} makes  Firms act, then wait to see how their action affected the world. Firms' knowledge is limited to their own revenue changes in order to see how they cope with the need to react to a simple price signal. They have only one action capable of effecting change and must work out if it has helped or hindered their aim to maximise revenue. Each Firm faces a world with one other Firm in it, both of them competing for a finite pot of demand. In the space between turns, their data may have been affected by factors other than their own action. 

It was shown that their ability to interpret the data correctly is dependent on the effect other agents have on the datum they are using, as well as the way in which \gls{lov} can affect the magnitude of those signals because of elasticity between goods. Interpretation of these signals was also shown to be dependent on noise levels. The \gls{dmr} model provides a coordination mechanism but illustrates the challenges of programming agents able to manage market signals. This weakness does allow for a close study of the problem, however: the \gls{dmr} issue can be examined in detail by testing in what way Firms succeed or fail.

Why can agents so successfully self-organise into spatial equilibrium arrangements when driven by density costs (section \ref{sec:results}), while even a simple \gls{dmr} problem (section \ref{sec:firmsLoV}) proves so challenging? Both are using minimal rationality, able to know only `better' from `worse'. Similarly, why can the transmission belt model successfully use price signals to achieve a stable equilibrium, despite large-scale agent interaction, when in the \gls{dmr} model, just two Firms struggle? The difference is in the density cost's combination with a demand function to rule out `wrong' decisions. This does not mean the density cost approach is necessarily better, if \gls{dmr} is the topic of interest.

What distinguishes the transmission belt mechanism from a \gls{dmr} approach? There is a clear difference in the speed of signalling; it is closer to the cybernetic ideal of `intrinsic control' mentioned in section \ref{sec:dmr}, adjusting prices with each and every demand signal received from the market. Each time a Person makes an exchange, using a \gls{ces} function, they exchange with every Firm in the model - so each Firm gets to respond using its stock targeting mechanism. A Firm's response is direct, in proportion to the change and guaranteed to be the correct polarity; \gls{dmr} Firms can make exactly the wrong interpretation and thus do exactly the wrong thing. 

This overview suggests a diversity of approaches to getting agents interacting, each suitable for different purposes. In-keeping with the theme of assessing \gls{abm} obstacles, Friedman's simple phrase seems apt: ``everything depends on the problem''. This certainly applies to the forms of interaction that may be relevant for a particular agent modelling problem. Spatial economics presents many interaction challenges; addressing them will almost certainly involve a range of interaction types and a willingness to try out new ideas.

\section{Assumptions \label{sec:concassumptions}}

\subsection{Overview}

Section \ref{sec:realism} discussed the common na\"{\i}ve `assumptions do not matter' interpretation of Friedman. As he says himself, of course they matter. In the case of the \gls{coremodel}, for example, one of the driving motivations behind its creation was an unhappiness with the key assumption that factor differences underlay all explanations for trade over distance, as well as problems with the set of assumptions and dependencies that made so much of economics ignore space altogether. Friedman was arguing instead that no perfect recipe exists for assessing the merits of a particular set of assumptions.

As chapter \ref{chapter:ABM} explored, the assumptions of neoclassical economics (and thus \gls{ge}) stand accused of creating ``citadels of crystalline mathematical perfection that would shatter if touched by the harsh rays of reality'' (\citealt{ball_social_2007} p.647). \gls{abm} has sold itself as the obvious answer, able to deal with the `messy' reality of real human behaviour (ibid) in a more `realistic' way. This thesis has suggested these criticisms may have been misplaced, and that agent modelling has fallen into just the `descriptive realism' trap that Friedman outlined so well in the 1950s. It has also argued this represents a major obstacle to pursuing further work in spatial economic agent modelling, but one that is surmountable with careful consideration the nature of assumptions and dependencies and a commitment to the `praxis' of model-building.

Avoiding `descriptive realism' does not mean avoiding reality. A balance needs to be struck between a dogmatic commitment to mapping model structures to real-world counterparts and a complete disregard for the connection between assumptions and the system being modelled. 

Bringing models to reality is often done through a basic requirement for structural parameters that have real-world counterparts, enabling direct model testing from data. So, for example, Bosker \emph{et al.} (\citeyear{bosker_adding_2010} p.811) seek out empirical counterparts for the \gls{coremodel}'s main variables. But the importance of theoretical work is in part as a source of new ideas on what sort of structural parameters might be important. This is the ultimate success of a good theory - making `substantive hypotheses' that propose new `ontological predictions' (section \ref{sec:prediction}) that suggest new places to look for real-world data, or new ways to examine existing information. 

%Section \ref{sec:ecocrit} quoted Solow: ``the role of theory is to explore what logic and simple assumptions can tell us about what data to look for and how to interpret them in connection with the question asked'' (\citealt{solow_georgescu-roegen_1997} p.267/8).
The density cost approach in section \ref{sec:results} identifies a difference between the effect of changing spatial and non-spatial costs, and provides a framework for considering the spatial and non-spatial components of good and people movement as identical in effect. This suggests some places to look for real-world data; this is discussed in section \ref{sec:theend}. Assuming all structural parameters to begin with can dampen this sort of discovery. Assumption of structural parameters is absolutely necessary for applied models, but potentially ossifying for theoretic models. 

\subsection{Model framework assumptions}
%\gls{abm} allows simplicity avoids contortions
Previous sections have discussed the thesis model framework assumptions and dependencies in close detail; this section looks at its more broad assumptions. Perhaps the most important is that it treats time as `just another factor', exchangeable for other factors, and supplied in a free daily flow. Agents need not concern themselves with the `time' that happens over many model iterations, except in the case of Firms waiting one timestep for a `sonar' from a previous action. Otherwise, all model effects emerge through actions agents take `now'. Section \ref{sec:spacecostgoods} discussed some of the key elements of real-world time dynamics this bypasses, while section \ref{sec:spacerisk} outlined the importance of risk and uncertainty in real-world time calculations. 

%The way problems with longer timespans is avoided is both powerful and brittle. 
The interaction and coordination problems described in section \ref{sec:coordoverview} are all deeply affected by time assumptions. Section \ref{sec:litbits} explained that Hagerstrandian diffusion dynamics were to be explicitly excluded from the thesis model framework. The simple time structure used has allowed a range of approaches to thinking about agent spatial economics. The issues raised by a very basic \gls{dmr} problem throw into sharp relief some of the challenges that might await once broader time-space considerations begin to come into play. This relates also to the use of Euclidean space; reality is networks and changing infrastructure (section \ref{sec:routedev}). Under assumptions of no network change and a high `route factor' it is possible to use Euclidean distance as a reasonably accurate proxy but it is still clearly unrealistic.
%And , transport network modelling leads straight onto considering diffusion processes.

The use of utility and maximisation is a key assumption. Chapter \ref{chapter:ABM} spent some time exploring critiques of utility. It also discussed how there was some confusion over the `level' that utility should be applied at, whether assumptions like the \gls{repag} could lead to problems, and whether the fact that people are not, actually, utility-maximisers means it should be avoided in agent models. The framework presented has demonstrated at the very least that rational spatial outcomes can result from the interaction of many agents given utility functions. The density cost would not work without this utility foundation allowing People to `trade off' differing spatial economic options; elasticity of substitution is key for the transmission belt model dynamic. It has been able to deal with the imposition of heterogeneity that spatial location introduces: agents in different places will face differing costs, and this will in part determine their optimal location choices. It has been argued that there is no need to appeal to psychological plausibility: to the extent that people will react to space cost changes consistently, treating agents as utility-maximisers is useful.

The particular utility approach used, the \gls{ces} function, proves the value of combining methods. It has been an indispensable `Swiss-army knife' for simplifying model dynamics: an `as if' proxy for much more complex processes. The consumption of $n$ goods through a \gls{ces} utility function is, of course, unrealistic. It subsumes a whole collection of processes that would, in reality, have to take place for such diverse bundles to be assembled. In the thesis model framework, People using the \gls{ces} do not need to travel to `fetch' each good; because of this, all route complexities are bypassed. Nor are People required to find a marketplace to buy them. They face no uncertainty in their space search: prices of all goods are known, regardless of distance. But in Feynman's terms, being in a framework that allows the \gls{ces} approach to be true does not mean being ignorant of the messy reality it is being asked to summarise. As a method for approaching the problem of space-costs using a relatively simple set of assumptions, it has proved invaluable.

The density cost approach was asked to stand in for the messy complexities of price-based land market interactions, as well as more intangible `congestion costs'. It meshes naturally with treating agents as separate decision-making units, but it has many features that make it unlike any real-world land market, reducing all complexities to a `second nature' interaction (section \ref{sec:1st2ndnature}) with other proximate agents. It actually does away with any market at all. All assumptions, however, will have some similarities to the phenonema being examined, and some contradictions. In this case, by removing several layers of complexity, a density cost allows the model to focus on its underlying rationale: examining how agents' spatial economic choices are affected by proximity costs; where they work, what they buy and where they choose to locate. 

The transmission belt model is also replete with `heroic assumptions'. In particular, no firm in reality targets a stock level like a hydraulic shock absorber, but this is the method that worked best for the model. It also illustrates the dangers of confusing emergence with explanation. As chapter \ref{chapter:ABM} argued, \gls{abm} has tended to privilege the `Hayekian' research agenda: to ``show how a solution is produced by the interactions of people each of whom possesses only partial knowledge'' (\citealt{hayek_use_1945}). The transmission belt model appears to do this. Can many agents, through price signals, coordinate to produce an optimal production level that maximises welfare? Yes. But this emergence by itself clearly does not explain what is causing it; this took a more detailed, simplified spreadsheet approach to track down a difference in magnitude of change (section \ref{sec:lov_vs_eos}).

In an agent context, the coordination question arises due to the nature of the modelling method - but that does not make it the default question to be asked. If there are ways round having to solve it, there may be good reasons to do so. It could be argued this is one of \gls{abm}'s built-in `streetlight effect' limitations: precisely that coordination needs so explicitly addressing (or avoiding, as with the density cost and strategy-free \gls{eos} method) before any research questions can be addressed. 
 
\section{Future work and concluding comments \label{sec:theend}}

Section \ref{sec:problem} framed some of the `big questions' of spatial economics in terms of contemporary arguments about relocalisation and changing costs. It ended by asking, ``is \gls{abm} well-positioned to deal with understanding how the `gestalt whole' spatial economy changes?''

The preceding analysis of dependencies and assumptions is chastening and promising in equal measure. Some narrow modelling problems have been solved, most notably the link between supply and demand in a spatial setting, including a dynamic production equilibrium. Keeping the core elements of the model framework simple has made it possible to clearly see some of the other obstacles in the way of progress. 

The most important lesson is perhaps this: it is highly unlikely that simply recreating spatial economic actors using an \gls{oop} framework will solve the problem. The challenges go much deeper than this. This thesis has argued that the lure of \gls{oop}-built virtual worlds able to furnish hermetic explanations has distracted from some hard-earned lessons earlier modellers discovered. The ideal of perfect `encapsulation of reality' has been tried before; arguments that \gls{abm} represents a `year zero' able to escape the problems this faced first time around do not stack up. 

That is the chastening part. There is promise, however: agent modelling is a powerful and flexible approach. As demonstrated, when compared to \gls{ge}, it has its own idiosyncratic problems, but it also provides a profound flexibility and is very amenable to a problem-solving, pragmatic approach.

This final section considers some possible future routes that spatial economic \gls{abm} might take. Many of these have already been considered during the discussion of dependencies and assumptions; a lot of scope lies here before adding any further complexity. A first problem might be: how to make an agent market structure that can improve on competition dynamics in a spatial setting? This thesis has - like \gls{ge} - tactically avoided issues of collusion and strategy. The transmission belt model, however, appeared to show genuine competition to be problematic. Is it something about the model's dependencies that causes this, or can this be addressed within the same framework? 

Competition in a market stucture is a fiendish problem from an \gls{abm} point of view: many things have to be explicitly specified. While the thesis has argued against always attempting to make every dynamic explicit, \gls{abm} does still encourage that approach - and nowhere more strongly than the birth and death of the agents themselves, if entering and exiting the market is vital to understanding competition. What kind of mechanism can be used?

As mentioned, all neoclassical economic theory (including \gls{ge}) bypasses the problem by assuming competition has taken place. As the \gls{coremodel} shows, this can still leave a great deal of complex algebraic organisation to be able to produce a workable model, but it still builds on static assumptions. Perhaps there are simplifications suited to \gls{abm}, as the density cost effectively simplified land and congestion costs. Or perhaps competition requires an \gls{abm} approach that tackles all its coordination problems head-on.

Some other problems arising from the listed dependencies include: is there a way of separating out Peoples' demand for wages from their good demands? Would that add anything useful to the analysis? Separating out different demand sectors requires some way to communicate costs; this was one of the key difficulties avoided by subsuming wages and goods into one demand function. What other methods for communicating costs can be used that would allow more flexibility while avoiding the `rivalrous money' / `currency entropy' problem (section \ref{sec:comm_costs})? Can the hill-climbing approach to utility keep demand economically sensible while allowing Firms a different approach to production that does not involve assuming demand is always met? Is solving \gls{dmr}-type coordination problems even necessary for understanding spatial economics from an \gls{abm} point of view - might it be fine to use Firms with uncanny knowledge abilities?

How do the spatial equilibrium outcomes of the density cost approach compare to a more explicit land market model? Is the `opposite' reaction of spatial and non-spatial costs (section \ref{sec:spatialvsnon}) something that only makes sense in the framework presented here or is it confirmed by other approaches? Is it detectable in reality? An increase in wages would lead to a relative decrease in commute costs, so the two costs are not as cleanly separable as in the density cost model, but the magnitude of change might be detectable.

Outside of the range of dependencies discussed above, more complex problems await. By far the most challenging of these is time. This is at the heart of the problem of modelling Isard's `gestalt whole' of agents who may be taking actions very far apart in time and space. It could be argued it is also vital for addressing the `problems' described in section \ref{sec:problem}. Keeping it simple has been very useful for the thesis models. But a rich array of possible modelling ideas exist beyond this approach to agents optimising `now' or at most between two timesteps. Any logical journey to make a model capable of examining this would need to make strategic simplifications of its own, but it may well be some approach using a \gls{ces} function could achieve that goal in a generic way. 
%In this thesis, time has been reduced to space, and all decisions forced to take place during one model iteration. The reality of time (discussed in chapter \ref{chapter:location}) is much more complex. 

As well as demand being distributed over `timespace', the issue of changes in production structure is deeply time-dependent. The flexibility demonstrated by the transmission belt model hints that perhaps \gls{abm} could be effective here. As costs change, this alters forward and backward production linkages (to use Marshall's language). The full problem of modelling this Jacobs-like `tangled bank' (\citealt{ellerman_jane_2002} p.6) of connections between productive units can be imagined as a mesh of differing input and output elasticities. Over distances, these processes overlap with diffusion issues as well as the network modelling concerns discussed in section \ref{sec:routedev}. 

Are there simplifications that can be made to keep this problem tractable when moving towards a more `tangled bank' model, or are more complex timing relations required? Is it possible to use input-output frameworks to model changes in morphology in this `tangled bank?' Is it possible to understand its relation to fundamental `primary sector' production and each following stage? The problem is in combining the fact that ``all products are simultaneously inputs to some organisations and outputs from others'' (\citealt{birkin_industrial_1986} p.176)  with how cost changes may affect the long-term location choices of agents, and the infrastructure development they must work with. 

Speaking of infrastructure, does any model of the `gestalt whole' have to take account of Haggett \emph{et al.}'s challenge: the need for a `comprehensive model of route development' (\citealt{haggett_locational_1977} p.95) capable of incorporating agents' response to cost changes, cognisant of the fact that it is `equally plausible to regard transport investment as a result of a need for movement or as a generator of movement'' (\citealt{chapman_people_1979} p.230, from section \ref{sec:routedev})? 

Chapter \ref{chapter:intro} mentioned the work of Kerschner and Hubacek (\citeyear{kerschner_assessing_2009}; section \ref{sec:problem}) examining the effect of oil cost change in an input-output framework but such approaches are bound both by rigid categories and an inability to conceptualise geographical change. This is something Jacobs was particularly critical of: in theorising how production morphology developed, she argued that:

\begin{quote}
``The point is that when new work is added to older work, the addition often cuts ruthlessly across categories of work, no matter how one may analyse the categories. Only in stagnant economies does work stay docilely within given categories'' (\citealt{jacobs_economy_1970} p.62).
\end{quote}
%She has a dig at economic planning, zoning, unions:

%Not mistaking map for territory
If she is right, the kind of categorisation methods used to collect data for input-output models will struggle with modelling change. While Jacobs' absolutism would appear to rule out data collection altogether, the kind of analysis carried out by Kerschner and Hubacek clearly helps to develop understanding of how different economic sectors are stressed by space cost change. But can generic methods be developed, parallel to the flexibility of the \gls{ces} function, able to model changes in input-output in a fluid way? Such an approach could look towards other generic approaches to product differentiation such as that developed by Lancaster (\citeyear{lancaster_socially_1975}), where goods are distinguished by their distance in something akin to a `characteristics space', rather than the generic, across-the-board \gls{lov} used in this thesis.
%Two promising facets that may involve more complex timing models are as follows. Firstly, there is scope for taking up Haggett \emph{et al.}'s challenge of further developing a `comprehensive model of route development' (\citealt{haggett_locational_1977} p.95) capable of incorporating agents' response to cost changes, and some generic way to understand how interests shape network outcomes. 

%Secondly, this thesis' simple approach to production and trade linkages is a possible next step. 
Food production and consumption is a fascinating case study in this respect: the shift in agricultural production is a quintessential `Engel's law' dynamic (\citealt{zimmerman_ernst_1932}; \citealt{murata_rural-urban_2002}; \citealt{murata_engels_2008}), where an increase in overall wealth leads to a drop in the percentage of revenue going into food production. In terms of production morphology, that dynamic masks a complex process where the changing opportunity costs of time have moved food processing out of the `productive unit' of the household (see e.g. \citealt{taylor_agricultural_2003}), slowly `outsourcing' more elements of food processing into the marketplace. The morphology of food production, then, has changed `from field to fork'. The problem is a difficult one, however, precisely because those changes are not isolated within any one, neatly definable sector.

The `transition' literature mentioned in the introduction has a large bias towards food: a burgeoning literature on `alternative food networks' has developed, in which by and large `alternative equals good / conventional equals bad' (\citealt{penker_mapping_2006} p.377). The underlying rationale for recommending localising these networks is a common-sense view that as increase in fuel prices (either due to peak oil or a carbon price) will make the primary need for food (and the localisation of other key non-food activities as far as possible) over-ride all other economic considerations. The economic case for this common-sense view is unproven; most extant theories say it is wrong. But equally, this thesis suggests there is no clear refutation of it either, if one considers space cost change in the round. 

Though many approaches exist able to say important things about aspects of the problem, given what has been discussed in this thesis, it is understandable that many economic geography researchers are skeptical of both economic and agent-based methods, if spatial economic approaches simply ignore problems beyond their `streetlight'. It is unlikely any one theoretical approach will offer definitive answers and as Isard says, ``formal theory, in and of itself, is highly unsatisfactory, too general and accordingly too sterile'' (\citealt{isard_location_1956} p.36). The tendency to disappear into virtual worlds, rather than find ways to use theory as an `engine of analysis' (section \ref{sec:virtualworlds}) must be fought. 

Agent modelling could be ideal: as this thesis has demonstrated, spatial economic agents are capable of producing robust rational outcomes in a flexible way. They could perhaps be particularly well-suited to the kind of change just described - not in creating complex, descriptive simulacra, but in using generic micro-economic ideas to study shifts in morphology.

The transmission belt models are examples of how economic change can produce a jump between discrete outcomes, mirroring the `catastrophic' shifts the \gls{coremodel} describes. This naturally presents the question: what sort of radical changes in morphology potentially exist in reality? The `transition' literature's take suggests cheap energy is a prerequisite for the existing industrial food network morphology. The extreme version of this is that we `eat fossil fuels' and as the cost of those goes up, that industrial system will not survive. Kerschner and Hubacek's examination of the data says exactly the opposite: the agricultural sector is one of the \emph{least} vulnerable to oil cost change. This lack of agreement suggests there is some way to go before a full understanding of changes in morphology can exist, but it also hints that modelling and data-driven analyses can offer answers quite at odds with common-sense empirical understandings.

%\gls{ge} attempts to endogenise all spatial outcomes into economic forces alone. If, in reality, the spatial economy is a complex mix of economic and other factors, this may suggest this `endogenise everything' philosophy will always produce incomplete models. 
A last note on modelling philosophy and how this informs research choices. As mentioned in section \ref{sec:marketstructnspace}, the neoclassical philosophy of `endogenise everything' may be faulty. The rationale for this philosophy is clear enough: the more can be shown to emerge from a model's internal dynamics, the more one can claim the model `explains' the dynamic it addresses. In the case of the transmission belt model, the `transmission belt plus density' model presented in section \ref{sec:tbmplusdensity} utilised its `general' ability to link supply and demand to endogenise the emergence of a difference in good price that had been externally imposed in section \ref{sec:LoVresult}. 

But the `endogenise everything' approach is predicated on the morphology of the spatial economy being entirely economic in origin. Section \ref{sec:marketstructnspace} noted Fujita musing that ``economic development may represent merely a possible result of the more fundamental process of social and cultural development''. It certainly may be true that while economic forces may have formed a given settlement pattern, non-economic forces keep it stable, even in the face of strong economic headwinds. There may be mix of dynamics, with economic forces powerfully shaping outcomes but not solely determining them. The difference between those two is crucial: it means pursuing some ultimate economic `theory of everything' to explain the spatial economy is likely to fail, as it excludes key factors.

%As Lucas says (section \ref{sec:ABMintro}), economics means `programming robot imitations of people'.
\gls{abm} has sold itself as a method for pursuing this brave new world of complex social forces, and although this thesis has argued it has got ahead of itself, as the discipline matures it should be well positioned to exploit these possibilities. It has been argued that no model should be asked to stand in as a `virtual world', only help as an engine of analysis. No modelling analysis will be able to `solve' the problem of optimal production given changing space costs. What modelling \emph{should} offer is a way to think through assumptions and develop new, informed intuitions. As Marshall put it:

\begin{quote}
``Economic laws and reasonings in fact are merely a part of the material of which conscience and common-sense have to make use in solving practical problems, and in laying down rules which may be a guide to life.'' (\citealt{marshall_principles_1895} p.x)
\end{quote}

%THE FOOKIN END.
